% ============================================================
% BAGIAN 2 - ANALISIS ALGORITMA PATHFINDING
% ============================================================

Bagian ini membahas analisis terhadap algoritma UCS, GBFS, dan A* yang digunakan pada penyelesaian \textit{Ice Sliding Puzzle Solver}. 
Analisis dilakukan dari sisi fungsi evaluasi, heuristik, optimalitas, kompleksitas, ruang state, serta kesesuaian algoritma terhadap karakteristik khusus puzzle.

\section{Definisi Cost dan Fungsi Evaluasi}

Pada permasalahan ini, cost tidak dihitung berdasarkan jumlah aksi saja, melainkan berdasarkan total cost tile yang dilewati selama aktor melakukan gerakan sliding. 
Tile awal sebelum gerakan tidak dihitung, sedangkan tile tempat aktor berhenti tetap dihitung karena tile tersebut termasuk tile yang dilewati.

\begin{tcolorbox}[
	colback=blue!5,
	colframe=blue!60!black,
	title=\textbf{Definisi Cost Gerakan},
	arc=2mm,
	boxrule=0.8pt
]
Jika dalam satu gerakan sliding aktor melewati tile dengan cost $c_1, c_2, \ldots, c_k$, maka cost gerakan tersebut adalah:
\[
	C_{\text{move}} = \sum_{i=1}^{k} c_i
\]
Total cost solusi adalah penjumlahan cost seluruh gerakan sliding:
\[
	C_{\text{total}} = \sum_{j=1}^{m} C_{\text{move}_j}
\]
dengan $m$ menyatakan banyaknya aksi gerakan dalam solusi.
\end{tcolorbox}

\begin{table}[H]
	\centering
	\caption{Definisi $g(n)$, $h(n)$, dan $f(n)$ pada setiap algoritma}
	\begin{tabular}{|p{2.7cm}|p{3.7cm}|p{3.7cm}|p{3.7cm}|}
		\hline
		\rowcolor{gray!25}
		\textbf{Algoritma} & \textbf{$g(n)$} & \textbf{$h(n)$} & \textbf{$f(n)$ / Prioritas} \\
		\hline
		UCS & Total cost aktual dari start ke state $n$. & Tidak digunakan. & $f(n)=g(n)$ \\
		\hline
		GBFS & Tetap dapat disimpan sebagai total cost solusi, tetapi tidak menjadi dasar prioritas. & Estimasi jarak atau cost dari state $n$ ke target berikutnya. & $f(n)=h(n)$ \\
		\hline
		A* & Total cost aktual dari start ke state $n$. & Estimasi jarak atau cost dari state $n$ ke target berikutnya. & $f(n)=g(n)+h(n)$ \\
		\hline
	\end{tabular}
\end{table}

\begin{tcolorbox}[
	colback=gray!5,
	colframe=black!60,
	title=\textbf{Perbedaan Utama Fungsi Evaluasi},
	arc=2mm,
	boxrule=0.8pt
]
\begin{itemize}
	\item UCS bersifat \textit{cost-based} karena hanya memilih state berdasarkan cost aktual terkecil.
	\item GBFS bersifat \textit{heuristic-based} karena hanya memilih state berdasarkan estimasi kedekatan ke target.
	\item A* bersifat gabungan karena mempertimbangkan cost aktual dan estimasi menuju target.
\end{itemize}
\end{tcolorbox}

\begin{table}[H]
	\centering
	\caption{Contoh konseptual pemilihan state berdasarkan fungsi evaluasi}
	\begin{tabular}{|p{2.3cm}|p{2.5cm}|p{2.5cm}|p{2.5cm}|p{3cm}|}
		\hline
		\rowcolor{blue!15}
		\textbf{State} & \textbf{$g(n)$} & \textbf{$h(n)$} & \textbf{$g(n)+h(n)$} & \textbf{Dipilih oleh} \\
		\hline
		A & 10 & 8 & 18 & UCS \\
		\hline
		B & 14 & 3 & 17 & GBFS dan A* \\
		\hline
		C & 12 & 6 & 18 & Tidak dipilih lebih awal \\
		\hline
	\end{tabular}
\end{table}

Pada contoh di atas, UCS memilih state A karena memiliki $g(n)$ terkecil. 
GBFS memilih state B karena memiliki $h(n)$ terkecil. 
A* juga memilih state B karena memiliki nilai $g(n)+h(n)$ terkecil.

\section{Analisis Heuristik}

Heuristik digunakan pada GBFS dan A* untuk memperkirakan kedekatan suatu state terhadap target berikutnya. 
Target berikutnya tidak selalu berupa titik tujuan \texttt{O}. 
Jika masih terdapat angka yang belum dikumpulkan, maka target berikutnya adalah angka yang harus dilewati sesuai urutan. 
Setelah seluruh angka terkumpul, barulah target heuristik diarahkan ke titik tujuan.

\begin{tcolorbox}[
	colback=orange!5,
	colframe=orange!70!black,
	title=\textbf{Target Heuristik},
	arc=2mm,
	boxrule=0.8pt
]
\[
	\text{Target}(n) =
	\begin{cases}
		\text{posisi angka berikutnya}, & \text{jika masih ada angka yang belum dikumpulkan} \\
		\text{posisi goal } O, & \text{jika semua angka sudah dikumpulkan}
	\end{cases}
\]
\end{tcolorbox}

\subsection{Heuristik Utama}

Heuristik utama yang umum digunakan pada grid adalah jarak Manhattan. 
Jarak Manhattan menghitung selisih absolut baris dan kolom antara posisi aktor dengan target.

\begin{tcolorbox}[
	colback=orange!5,
	colframe=orange!70!black,
	title=\textbf{Heuristik Manhattan Distance},
	arc=2mm,
	boxrule=0.8pt
]
Misalkan posisi aktor adalah $(r_1,c_1)$ dan posisi target adalah $(r_2,c_2)$, maka:
\[
	h(n)=|r_1-r_2|+|c_1-c_2|
\]
Heuristik ini dipilih karena sederhana, cepat dihitung, dan sesuai dengan papan berbentuk grid yang gerakannya hanya horizontal atau vertikal.
\end{tcolorbox}

\subsection{Heuristik Alternatif}

Apabila mengerjakan bonus heuristik, beberapa alternatif heuristik dapat digunakan untuk membandingkan performa pencarian. 
Setiap heuristik memiliki karakteristik yang berbeda terhadap akurasi estimasi dan waktu komputasi.

\begin{longtable}{|p{3.5cm}|p{5cm}|p{5.5cm}|}
	\caption{Contoh heuristik utama dan alternatif} \\
	\hline
	\rowcolor{orange!20}
	\textbf{Heuristik} & \textbf{Rumus / Ide} & \textbf{Karakteristik} \\
	\hline
	\endfirsthead
	
	\hline
	\rowcolor{orange!20}
	\textbf{Heuristik} & \textbf{Rumus / Ide} & \textbf{Karakteristik} \\
	\hline
	\endhead
	
	Manhattan Distance & $|r_1-r_2|+|c_1-c_2|$ & Cepat dihitung dan cocok untuk grid 4 arah, tetapi belum memperhitungkan sliding, obstacle, lava, dan cost tile. \\
	\hline
	Euclidean Distance & $\sqrt{(r_1-r_2)^2+(c_1-c_2)^2}$ & Mengukur jarak garis lurus, tetapi kurang merepresentasikan gerakan grid yang hanya horizontal dan vertikal. \\
	\hline
	Minimum Tile Cost Weighted Manhattan & Manhattan Distance dikalikan cost tile minimum yang dapat dilewati. & Lebih dekat ke estimasi cost karena mempertimbangkan batas bawah cost tile, tetapi tetap belum sepenuhnya memperhitungkan sliding dan obstacle. \\
	\hline
	Remaining Target Estimate & Estimasi dari posisi saat ini ke angka berikutnya, lalu dilanjutkan secara kasar ke target-target berikutnya. & Dapat lebih informatif, tetapi perlu hati-hati agar tidak melebih-lebihkan estimasi jika digunakan pada A*. \\
	\hline
\end{longtable}

\subsection{Aspek yang Dipertimbangkan Heuristik}

Tidak semua heuristik memperhitungkan seluruh aturan pada Ice Sliding Puzzle. 
Beberapa heuristik hanya memperkirakan jarak geometris, sedangkan validitas gerakan tetap ditentukan oleh simulasi sliding.

\begin{table}[H]
	\centering
	\caption{Aspek puzzle yang dipertimbangkan oleh heuristik}
	\begin{tabular}{|p{5cm}|p{2.5cm}|p{6cm}|}
		\hline
		\rowcolor{orange!20}
		\textbf{Aspek} & \textbf{Dipertimbangkan?} & \textbf{Keterangan} \\
		\hline
		Posisi goal & Ya & Goal digunakan sebagai target apabila seluruh angka sudah terkumpul. \\
		\hline
		Posisi angka berikutnya & Ya & Jika masih ada angka, heuristik diarahkan ke angka berikutnya. \\
		\hline
		Sliding movement & Tidak langsung & Heuristik jarak hanya memberi estimasi; gerakan sliding tetap divalidasi oleh mekanisme ekspansi. \\
		\hline
		Cost tile & Tergantung heuristik & Manhattan biasa tidak memperhitungkan cost tile, sedangkan weighted heuristic dapat memasukkan cost minimum. \\
		\hline
		Obstacle & Tidak langsung & Obstacle tidak dihitung dalam rumus jarak sederhana, tetapi memengaruhi hasil validasi gerakan. \\
		\hline
		Lava & Tidak langsung & Lava tidak dihitung dalam heuristik sederhana, tetapi state yang melewati lava akan dibuang saat ekspansi. \\
		\hline
	\end{tabular}
\end{table}

\subsection{Admissibility Heuristik A*}

Heuristik dikatakan \textit{admissible} apabila tidak pernah melebih-lebihkan cost minimum sebenarnya dari suatu state menuju goal. 
Dengan kata lain, untuk setiap state $n$ berlaku:

\[
	h(n) \leq h^*(n)
\]

dengan $h^*(n)$ adalah cost optimal sebenarnya dari state $n$ menuju target.

\begin{tcolorbox}[
	colback=green!5,
	colframe=green!50!black,
	title=\textbf{Analisis Admissibility},
	arc=2mm,
	boxrule=0.8pt
]
Jika heuristik yang digunakan adalah Manhattan Distance murni, maka heuristik tersebut aman sebagai estimasi jarak geometris, tetapi belum tentu selalu menjadi estimasi cost yang ketat pada puzzle sliding. 
Apabila semua cost tile minimal bernilai $1$, Manhattan Distance dapat dipandang sebagai batas bawah sederhana terhadap jumlah perpindahan grid. 
Namun, karena satu aksi sliding dapat melewati banyak tile sekaligus dan struktur obstacle dapat membuat aktor harus mengambil rute tidak langsung, admissibility perlu dianalisis berdasarkan definisi cost yang digunakan.
\end{tcolorbox}

Untuk menjaga heuristik A* tetap admissible terhadap cost, heuristik yang lebih aman adalah menggunakan estimasi berbasis batas bawah cost tile, misalnya:

\[
	h(n)=d_{\text{Manhattan}}(n,\text{target}) \times c_{\min}
\]

dengan $c_{\min}$ adalah cost minimum dari tile yang dapat dilewati. 
Jika $c_{\min}$ tidak melebihi cost traversal mana pun, maka estimasi ini tidak akan melebihi cost aktual minimum secara langsung terhadap jarak grid. 
Namun, tetap perlu dicatat bahwa karakter sliding dapat membuat estimasi ini tidak selalu sangat informatif.

\subsection{Konsistensi Heuristik}

Heuristik dikatakan konsisten atau monotonic apabila untuk setiap state $n$ dan successor $n'$ berlaku:

\[
	h(n) \leq c(n,n') + h(n')
\]

dengan $c(n,n')$ adalah cost gerakan dari $n$ ke $n'$.

\begin{table}[H]
	\centering
	\caption{Hubungan heuristik dengan performa pencarian}
	\begin{tabular}{|p{4cm}|p{10cm}|}
		\hline
		\rowcolor{orange!20}
		\textbf{Kualitas Heuristik} & \textbf{Dampak terhadap Algoritma} \\
		\hline
		Terlalu kecil & A* menjadi mirip UCS karena estimasi kurang membantu mengarahkan pencarian. \\
		\hline
		Terlalu besar & A* berisiko kehilangan jaminan optimalitas jika heuristik tidak admissible. \\
		\hline
		Informatif dan admissible & A* dapat mengurangi jumlah iterasi tanpa mengorbankan optimalitas. \\
		\hline
		Tidak mempertimbangkan obstacle/lava & Algoritma dapat mengejar target yang tampak dekat, tetapi sebenarnya sulit atau tidak valid dicapai. \\
		\hline
	\end{tabular}
\end{table}

\section{Analisis Optimalitas}

Optimalitas menunjukkan apakah algoritma dapat menjamin solusi dengan total cost paling minimum. 
Pada Ice Sliding Puzzle, optimalitas sangat bergantung pada cara algoritma memilih state dan bagaimana cost tile diperhitungkan.

\begin{table}[H]
	\centering
	\caption{Analisis optimalitas algoritma}
	\begin{tabular}{|p{3cm}|p{3cm}|p{8cm}|}
		\hline
		\rowcolor{green!15}
		\textbf{Algoritma} & \textbf{Optimal?} & \textbf{Alasan} \\
		\hline
		UCS & Ya & UCS selalu memilih state dengan total cost aktual terkecil. Dengan cost tile positif, solusi pertama yang mencapai goal merupakan solusi dengan cost minimum. \\
		\hline
		GBFS & Tidak selalu & GBFS hanya memilih state berdasarkan heuristik terkecil. Jalur yang terlihat dekat ke target belum tentu memiliki total cost paling kecil. \\
		\hline
		A* & Ya, dengan syarat & A* optimal apabila heuristik yang digunakan admissible dan cost gerakan tidak negatif. Jika heuristik melebih-lebihkan cost, jaminan optimalitas dapat hilang. \\
		\hline
	\end{tabular}
\end{table}

\begin{tcolorbox}[
	colback=green!5,
	colframe=green!50!black,
	title=\textbf{Hubungan Admissibility dan Optimalitas A*},
	arc=2mm,
	boxrule=0.8pt
]
A* dapat menjamin solusi optimal jika heuristik tidak pernah melebih-lebihkan cost sebenarnya menuju goal. 
Jika $h(n)$ admissible, maka nilai $f(n)=g(n)+h(n)$ tidak akan membuat algoritma mengabaikan jalur yang sebenarnya lebih murah.
\end{tcolorbox}

Pengaruh cost tile juga penting dalam optimalitas. 
Pada papan dengan cost berbeda-beda, jalur yang secara visual lebih pendek dapat memiliki total cost lebih besar dibanding jalur yang lebih panjang tetapi melewati tile murah. 
Karena itu, algoritma yang memperhatikan $g(n)$, seperti UCS dan A*, lebih aman digunakan untuk mencari cost minimum dibanding GBFS.

\section{Analisis Kompleksitas}

Kompleksitas algoritma dipengaruhi oleh banyaknya state yang mungkin dibentuk, jumlah arah gerakan, serta struktur papan. 
Misalkan papan berukuran $N \times M$, jumlah posisi maksimum adalah $N \cdot M$. 
Jika terdapat $K$ angka wajib, maka progres angka dapat direpresentasikan dengan bitmask hingga $2^K$ kemungkinan. 
Namun, karena angka harus dikumpulkan secara berurutan, jumlah progres valid dalam praktik lebih kecil dibanding semua kombinasi bitmask.

\begin{tcolorbox}[
	colback=purple!5,
	colframe=purple!60!black,
	title=\textbf{Perkiraan Jumlah State},
	arc=2mm,
	boxrule=0.8pt
]
Secara umum, batas atas jumlah state dapat ditulis sebagai:
\[
	|S| = O(N \cdot M \cdot 2^K)
\]
dengan:
\begin{itemize}
	\item $N \cdot M$ adalah jumlah maksimum posisi pada papan,
	\item $K$ adalah jumlah angka wajib yang harus dikumpulkan,
	\item $2^K$ adalah kemungkinan kombinasi angka yang sudah dikumpulkan.
\end{itemize}
Karena angka harus dikumpulkan berurutan, jumlah state valid biasanya lebih kecil daripada batas atas tersebut.
\end{tcolorbox}

Pada setiap state, algoritma mencoba empat arah gerakan. 
Namun, satu gerakan sliding dapat memindai beberapa tile sampai berhenti. 
Dalam kasus terburuk, satu sliding dapat memeriksa hingga $O(\max(N,M))$ tile.

\[
	T_{\text{expand}} = O(4 \cdot \max(N,M))
\]

\begin{longtable}{|p{3cm}|p{5.2cm}|p{5.7cm}|}
	\caption{Kompleksitas waktu dan ruang algoritma} \\
	\hline
	\rowcolor{purple!15}
	\textbf{Algoritma} & \textbf{Kompleksitas Waktu} & \textbf{Kompleksitas Ruang} \\
	\hline
	\endfirsthead
	
	\hline
	\rowcolor{purple!15}
	\textbf{Algoritma} & \textbf{Kompleksitas Waktu} & \textbf{Kompleksitas Ruang} \\
	\hline
	\endhead
	
	UCS & $O(|S| \log |S| \cdot 4 \cdot \max(N,M))$ & $O(|S|)$ untuk priority queue dan best cost map. \\
	\hline
	GBFS & $O(|S| \log |S| \cdot 4 \cdot \max(N,M))$ pada kasus terburuk, tetapi sering lebih cepat jika heuristik efektif. & $O(|S|)$ untuk priority queue dan struktur visited/best cost. \\
	\hline
	A* & $O(|S| \log |S| \cdot 4 \cdot \max(N,M))$ pada kasus terburuk, tetapi dapat lebih efisien dari UCS jika heuristik informatif. & $O(|S|)$ untuk priority queue dan best cost map. \\
	\hline
\end{longtable}

\begin{table}[H]
	\centering
	\caption{Faktor yang memengaruhi kompleksitas}
	\begin{tabular}{|p{4.5cm}|p{9.5cm}|}
		\hline
		\rowcolor{purple!15}
		\textbf{Faktor} & \textbf{Pengaruh} \\
		\hline
		Ukuran papan $N \times M$ & Semakin besar papan, semakin banyak posisi yang mungkin menjadi state. \\
		\hline
		Jumlah angka wajib $K$ & Semakin banyak angka, semakin banyak variasi progres angka yang perlu disimpan. \\
		\hline
		Jumlah arah gerakan & Setiap state maksimal diekspansi ke 4 arah. \\
		\hline
		Panjang sliding & Semakin panjang lintasan sliding, semakin besar biaya simulasi satu aksi. \\
		\hline
		Obstacle & Dapat mengurangi area yang dapat dicapai, tetapi juga dapat membuat rute menjadi lebih kompleks. \\
		\hline
		Lava & Mengurangi jalur valid karena state yang melewati lava langsung dibuang. \\
		\hline
		Kualitas heuristik & Sangat memengaruhi GBFS dan A*, terutama dalam mengurangi jumlah iterasi. \\
		\hline
	\end{tabular}
\end{table}

\section{Analisis State Space}

State space adalah himpunan seluruh konfigurasi yang mungkin diproses oleh algoritma. 
Pada puzzle ini, state tidak cukup hanya menyimpan posisi aktor, karena aktor yang berada di posisi sama dapat memiliki kondisi berbeda tergantung angka yang sudah dikumpulkan.

\begin{tcolorbox}[
	colback=blue!5,
	colframe=blue!60!black,
	title=\textbf{Definisi State Space},
	arc=2mm,
	boxrule=0.8pt
]
Sebuah state dapat ditulis sebagai:
\[
	s = (pos, mask, last, path, g, h)
\]
dengan $pos$ adalah posisi aktor, $mask$ adalah penanda angka yang sudah dikumpulkan, $last$ adalah angka terakhir yang valid, $path$ adalah urutan gerakan, $g$ adalah total cost aktual, dan $h$ adalah nilai heuristik.
\end{tcolorbox}

\begin{table}[H]
	\centering
	\caption{Analisis komponen state space}
	\begin{tabular}{|p{4cm}|p{10cm}|}
		\hline
		\rowcolor{blue!15}
		\textbf{Komponen} & \textbf{Analisis} \\
		\hline
		Jumlah posisi maksimum & Sebanyak $N \cdot M$, tetapi posisi berupa obstacle tidak dapat ditempati. \\
		\hline
		Constraint angka & Membuat state harus membedakan progres angka yang sudah dikumpulkan. \\
		\hline
		Collected mask & Digunakan untuk menandai angka yang sudah dilewati tanpa perlu menyimpan list angka secara eksplisit. \\
		\hline
		State yang sama & Dua state dianggap sama jika posisi aktor dan collected mask sama. \\
		\hline
		State berbeda & Dua state dengan posisi sama tetapi collected mask berbeda tetap dianggap berbeda. \\
		\hline
		Pruning & State diabaikan jika sudah pernah dicapai dengan cost yang lebih kecil atau sama. \\
		\hline
	\end{tabular}
\end{table}

\begin{tcolorbox}[
	colback=gray!5,
	colframe=black!60,
	title=\textbf{Mengapa Posisi Saja Tidak Cukup?},
	arc=2mm,
	boxrule=0.8pt
]
Misalkan aktor berada pada posisi $(r,c)$. 
Jika pada kondisi pertama aktor sudah mengumpulkan angka \texttt{0}, sedangkan pada kondisi kedua belum, maka kedua kondisi tersebut tidak boleh dianggap sama. 
Hal ini karena langkah valid berikutnya berbeda: state pertama dapat mencari angka \texttt{1}, sedangkan state kedua masih harus mencari angka \texttt{0}.
\end{tcolorbox}

\section{Perbandingan Teoritis Algoritma}

Secara teoritis, UCS, GBFS, dan A* memiliki keunggulan pada aspek yang berbeda. 
UCS unggul dalam jaminan optimalitas, GBFS unggul dalam kecepatan pencarian jika heuristik baik, sedangkan A* mencoba menyeimbangkan keduanya.

\subsection{UCS vs GBFS}

\begin{table}[H]
	\centering
	\caption{Perbandingan UCS dan GBFS}
	\begin{tabular}{|p{4cm}|p{5cm}|p{5cm}|}
		\hline
		\rowcolor{gray!25}
		\textbf{Aspek} & \textbf{UCS} & \textbf{GBFS} \\
		\hline
		Dasar pemilihan state & Cost aktual $g(n)$ & Heuristik $h(n)$ \\
		\hline
		Optimalitas & Menjamin solusi cost minimum & Tidak menjamin solusi cost minimum \\
		\hline
		Jumlah iterasi & Cenderung lebih banyak & Cenderung lebih sedikit jika heuristik baik \\
		\hline
		Ketergantungan heuristik & Tidak bergantung heuristik & Sangat bergantung heuristik \\
		\hline
		Kesesuaian untuk cost berbeda & Sangat sesuai & Kurang aman untuk cost minimum \\
		\hline
	\end{tabular}
\end{table}

\subsection{UCS vs A*}

\begin{table}[H]
	\centering
	\caption{Perbandingan UCS dan A*}
	\begin{tabular}{|p{4cm}|p{5cm}|p{5cm}|}
		\hline
		\rowcolor{gray!25}
		\textbf{Aspek} & \textbf{UCS} & \textbf{A*} \\
		\hline
		Dasar pemilihan state & $g(n)$ & $g(n)+h(n)$ \\
		\hline
		Optimalitas & Optimal untuk cost positif & Optimal jika heuristik admissible \\
		\hline
		Arah pencarian & Tidak terarah ke goal & Lebih terarah ke target \\
		\hline
		Jumlah iterasi & Dapat sangat besar & Umumnya lebih sedikit jika heuristik informatif \\
		\hline
		Kelemahan & Tidak menggunakan estimasi & Bergantung kualitas heuristik \\
		\hline
	\end{tabular}
\end{table}

\subsection{GBFS vs A*}

\begin{table}[H]
	\centering
	\caption{Perbandingan GBFS dan A*}
	\begin{tabular}{|p{4cm}|p{5cm}|p{5cm}|}
		\hline
		\rowcolor{gray!25}
		\textbf{Aspek} & \textbf{GBFS} & \textbf{A*} \\
		\hline
		Dasar pemilihan state & $h(n)$ & $g(n)+h(n)$ \\
		\hline
		Kecepatan & Bisa sangat cepat & Biasanya lebih stabil \\
		\hline
		Optimalitas & Tidak dijamin & Dijamin jika heuristik admissible \\
		\hline
		Pengaruh cost aktual & Tidak menjadi prioritas utama & Selalu dipertimbangkan \\
		\hline
		Risiko & Terjebak pada jalur yang tampak dekat tetapi mahal & Lebih kecil karena mempertimbangkan cost aktual \\
		\hline
	\end{tabular}
\end{table}

\begin{longtable}{|p{4cm}|p{3.2cm}|p{3.2cm}|p{3.2cm}|}
	\caption{Ringkasan perbandingan teoritis UCS, GBFS, dan A*} \\
	\hline
	\rowcolor{gray!25}
	\textbf{Aspek} & \textbf{UCS} & \textbf{GBFS} & \textbf{A*} \\
	\hline
	\endfirsthead
	
	\hline
	\rowcolor{gray!25}
	\textbf{Aspek} & \textbf{UCS} & \textbf{GBFS} & \textbf{A*} \\
	\hline
	\endhead
	
	Optimalitas & Tinggi & Rendah & Tinggi jika heuristik admissible \\
	\hline
	Jumlah iterasi & Besar & Kecil jika heuristik baik & Sedang hingga kecil \\
	\hline
	Waktu eksekusi & Cenderung lebih lama & Cenderung lebih cepat & Bergantung heuristik \\
	\hline
	Penggunaan memori & Besar & Besar pada kasus buruk & Besar pada kasus buruk \\
	\hline
	Ketergantungan heuristik & Tidak ada & Sangat tinggi & Tinggi \\
	\hline
	Cocok untuk cost berbeda & Sangat cocok & Kurang cocok untuk cost minimum & Cocok jika heuristik baik \\
	\hline
\end{longtable}

\begin{tcolorbox}[
	colback=green!5,
	colframe=green!50!black,
	title=\textbf{Kesimpulan Teoritis},
	arc=2mm,
	boxrule=0.8pt
]
\begin{itemize}
	\item Algoritma yang paling aman untuk mencari cost minimum adalah UCS.
	\item Algoritma yang secara teoritis paling cepat menemukan solusi adalah GBFS, tetapi tidak menjamin cost minimum.
	\item Algoritma yang paling seimbang untuk Ice Sliding Puzzle adalah A*, karena menggabungkan cost aktual dan heuristik.
\end{itemize}
\end{tcolorbox}

\section{Analisis Khusus Ice Sliding Puzzle}

Ice Sliding Puzzle memiliki karakteristik yang berbeda dari pathfinding biasa. 
Pada pathfinding biasa, satu aksi umumnya berarti berpindah ke satu tile tetangga. 
Pada puzzle ini, satu aksi dapat melewati banyak tile karena aktor terus bergerak sampai berhenti di depan rintangan.

\begin{table}[H]
	\centering
	\caption{Perbedaan pathfinding biasa dan Ice Sliding Puzzle}
	\begin{tabular}{|p{4.5cm}|p{4.5cm}|p{5cm}|}
		\hline
		\rowcolor{cyan!15}
		\textbf{Aspek} & \textbf{Pathfinding Biasa} & \textbf{Ice Sliding Puzzle} \\
		\hline
		Satu aksi & Berpindah satu tile & Sliding hingga berhenti \\
		\hline
		Cost gerakan & Biasanya cost tile tujuan atau cost edge & Jumlah cost seluruh tile yang dilewati \\
		\hline
		Obstacle & Menghalangi perpindahan & Menjadi titik penghenti sliding \\
		\hline
		Lava & Tergantung aturan kasus & Selalu membuat path tidak valid jika dilewati \\
		\hline
		Keluar papan & Biasanya tidak diperbolehkan & Menyebabkan game over \\
		\hline
		Goal & Biasanya cukup mencapai posisi goal & Harus berhenti tepat di \texttt{O} \\
		\hline
		Constraint tambahan & Umumnya tidak ada & Angka harus dilewati berurutan \\
		\hline
	\end{tabular}
\end{table}

\begin{tcolorbox}[
	colback=cyan!5,
	colframe=cyan!60!black,
	title=\textbf{Dampak Karakteristik Sliding terhadap Algoritma},
	arc=2mm,
	boxrule=0.8pt
]
Karena satu aksi dapat melewati banyak tile, algoritma tidak dapat langsung memperlakukan tetangga state sebagai empat tile di sekitar aktor. 
Setiap ekspansi harus melakukan simulasi gerakan sampai aktor berhenti atau gerakan menjadi tidak valid. 
Hal ini membuat proses ekspansi lebih mahal dibanding pathfinding grid biasa.
\end{tcolorbox}

\begin{table}[H]
	\centering
	\caption{Aturan khusus puzzle dan dampaknya terhadap pencarian}
	\begin{tabular}{|p{4cm}|p{10cm}|}
		\hline
		\rowcolor{cyan!15}
		\textbf{Aturan} & \textbf{Dampak terhadap Algoritma} \\
		\hline
		Satu aksi adalah sliding & Successor state hanya terbentuk pada posisi akhir sliding, bukan pada setiap tile yang dilewati. \\
		\hline
		Cost seluruh tile dilewati & Algoritma harus menjumlahkan cost sepanjang lintasan sliding. \\
		\hline
		Lava menyebabkan game over & State hasil gerakan yang melewati lava harus langsung dibuang. \\
		\hline
		Keluar papan menyebabkan game over & Gerakan tanpa rintangan penahan di tepi papan dianggap tidak valid. \\
		\hline
		Angka berurutan & State harus menyimpan progres angka yang sudah dikumpulkan. \\
		\hline
		Goal harus berhenti tepat di \texttt{O} & Melewati \texttt{O} saja tidak cukup; posisi akhir sliding harus sama dengan posisi goal. \\
		\hline
	\end{tabular}
\end{table}

\begin{tcolorbox}[
	colback=blue!5,
	colframe=blue!60!black,
	title=\textbf{Implikasi Akhir terhadap Pemilihan Algoritma},
	arc=2mm,
	boxrule=0.8pt
]
Ice Sliding Puzzle lebih cocok dianalisis sebagai pencarian pada ruang state berbobot, bukan sekadar pencarian jalur grid biasa. 
Karena terdapat cost tile yang berbeda-beda, UCS dan A* lebih aman digunakan untuk mencari solusi dengan total cost rendah. 
GBFS tetap berguna sebagai pembanding karena dapat menunjukkan bagaimana heuristik memengaruhi kecepatan pencarian, meskipun tidak selalu menghasilkan solusi optimal.
\end{tcolorbox}